The non-strict threshold in Siu's theorem is essential, even on a compact K\"ahler manifold.

\begin{example}\label{ex:strict-lelong-level}\index{Lelong number!strict level sets}
Let $X=\mathbb P^2$, with homogeneous coordinates $[x:y:z]$, and let $\sigma$ be the Fubini--Study form normalized by $\int_X\sigma^2=1$. Put $\|Z\|^2=|x|^2+|y|^2+|z|^2$, $a_j=2^{-j}$ and $\lambda_j=2^{-j-1}$ for $j\geq1$. The global function
\[
 u([Z])=\log\frac{|y|^2}{\|Z\|^2}
       +\sum_{j\geq1}\lambda_j
          \log\frac{|x-a_jz|^2}{(1+a_j^2)\|Z\|^2}
\]
belongs to $\PSH(X,\frac32\sigma)$. For the line $D=\{y=0\}$ and the points $p_j=[a_j:0:1]$, it satisfies
\[
 \{\nu(u,\cdot)\geq1\}=D,
 \qquad \{\nu(u,\cdot)>1\}=\{p_j:j\geq1\}.
\]
The second set is not analytic.
\end{example}
\begin{proof}
For a unit linear form $\ell$, the function $v_\ell=\log(|\ell(Z)|^2/\|Z\|^2)$ is nonpositive and satisfies $\ddc v_\ell=[\ell=0]-\sigma$. Its $L^1(X,\sigma^2)$ norm is finite and independent of $\ell$: the group $U(3)$ acts on homogeneous coordinates, preserves $\sigma$, and acts transitively on unit linear forms. Since $\sum_j\lambda_j=1/2$, the displayed partial sums decrease to an $L^1$ function. Their common curvature lower bound is $-\frac32\sigma$, so their limit is quasi-plurisubharmonic and
\[
 T:=\tfrac32\sigma+\ddc u=[D]+\sum_{j\geq1}\lambda_j[L_j],
 \qquad L_j=\{x-a_jz=0\}.
\]
The series of positive closed currents converges: every line has $\sigma$-trace mass one and the coefficients are summable.

We compute the Lelong numbers from the local mass formula. In affine coordinates $w=(w_1,w_2)$ centered at a point $q$, put $\kappa=\ddc|w|^2$. An affine complex line $H$ at distance $d$ from $q$ satisfies
\[
 \int_{H\cap B(q,r)}\kappa=(r^2-d^2)_+.
\]
Consequently $r^{-2}\int_{H\cap B(q,r)}\kappa$ is at most one and tends to $\mathbf1_H(q)$ as $r\downarrow0$; a line outside the chart contributes zero. Tonelli's theorem and dominated convergence for the summable coefficients therefore give
\[
 \nu(u,q)=\nu(T,q)
       =\mathbf1_D(q)+\sum_{j\geq1}\lambda_j\mathbf1_{L_j}(q).
\]
Here a smooth local correction does not change the Lelong number. This computation includes points where infinitely many lines accumulate.

The lines $L_j$ meet $D$ at the distinct points $p_j$ and meet one another only at $p_\infty=[0:1:0]\notin D$. Thus $\nu(u,p_j)=1+\lambda_j$, whereas $\nu(u,q)=1$ for $q\in D\setminus\{p_j:j\geq1\}$. Outside $D$, the Lelong number is at most $\sum_j\lambda_j=1/2$. The two level-set identities follow. Finally, $p_j\to p_*=[0:0:1]$, and $\nu(u,p_*)=1$. The strict level set is not closed and hence is not analytic.
\end{proof}
